\chapter{State-of-the-Art}
\label{chap:sota}

The State-of-the-Art section presents the papers that lay at the foundation at this research work, and positions it adverseraly. The sections in this chapter range will provide an overview over the communication mediums that we considered for this work: Bluetooth Classic, Bluetooth Low Energy, Bluetooth Mesh, Wi-Fi ad-hoc and Direct, and others. Next, this section continues with elaborating the networking assumptions that make this work possible, primarily discussing the difference between traditional packet-switching and the Store-Carry-Forward paradigm, that enables Delay-Tolerant Networking, while also presenting similar systems that have been developed using this paradigm. Followingly, we present main approaches that enable packet propagation in Delay Tolerant Networks (DTNs), explain each one of them, and show which one we used. Then, we present leading Decentralized Social Networks, and discuss whether they could have been used in our scenarios, showing that they actually cannot, and finally present the pinnacle: Grassroots Social Networks. Grassroots Social Networks lay at the core of this research work, and the networking layer presented here is a core work that lays at its foundation.

\section{Proximate Peer-to-peer Communication Technologies}
\label{sec:proximate-p2p-protocols}
To create the peer-to-peer communication platform, different technologies were evaluated based on the following criteria:
\begin{itemize}
    \item \textbf{C1} Spread of the protocol: As a user-facing tool, the communication medium had to be present in virtually all smartphones used today; thus, technologies that are supported by a handful of flagship smartphones, while being an interesting point for future work, are out of the question for this work.
    \item \textbf{C2} Low power: the communication protocol should not use a high amount of power to propagate packets
    \item \textbf{C3} Target: the communication protocol should be able to carry text messages efficiently, so a text message should be received within a short time in close range (e.g. 200KB in 1 second). 
    \item \textbf{C4} Mobility: protocols neede to support a high node mobility stemming from the mobility patterns of the users
    \item \textbf{C5} Infrastructure: there can be no external infrastructure necessary for the normal operation of the technology
    \item \textbf{C6} The protocol should be operational even when the app is backgrounded
\end{itemize}
\subsection{Wi-Fi}
To relay information directly between peer smartphones, the omnipresent Wi-Fi protocol was first investigated. Propelled by the Wi-Fi Alliance, WLAN is the wireless 802.11 protocol which enables most of the world's computers to communicate. Together with cellular-based mobile networks like 5G, home and commercial routers provide wireless radio access networks for mobile devices to connect to the Internet backbone through a multitude of devices: antennas (Radio Units), computers for signal processing (Distributed Units), base stations (Centralized Units), and more~\cite{5g}.
For the standard Internet-backbone to function, all the aforementioned infrastructure units need to have a constant stream of power, an unblocked radio spectrum, and operational staff. In emergency scenarios such as natural disasters, wars, or politically-motivated violence against the freedom of speech through censorship, media control, or monitoring, the aforementioned infrastructure might not be available for usage.
As alternatives to the classical Wi-Fi WLAN technology to facilitate DTNs, we evaluated the following p2p communication media based on the aforementioned common criteria: Wifi Adhoc, Wifi Direct, Wifi Aware, Bluetooth BDR, Bluetooth Low Energy.
The Wi-Fi Ad-Hoc mode was the first design to adapt the WLAN standard to direct peer-to-peer wireless communication.
However, it suffered from several notable architectural shortcomings, including inefficient power management, a lack of robust security mechanisms, and the requirement for devices to manually decide whether to act as the access point or a client~\cite{2012IEEE80211}.
In the Wi-Fi Ad-Hoc mode, idle smartphones suffer from battery drain due to a lack of distributed power saving mode~\cite{Friedman2013}. However, actively communicating smartphones had the same power cost as using a standard access point, and comparing to Bluetooth, Wi-Fi Ad-Hoc still has a vastly superior throughput rate, which is optimal for larger file transfers.
In our context it is still not as useful due to \textbf{C3}: the transfer size is minimal, and we optimize for efficient battery usage through \textbf{C2}.
Wi-Fi Ad-Hoc was eventually replaced by Wi-Fi Direct, which modernized peer-to-peer connectivity by adding better security and an autonomous negotiation protocol that allows connecting devices to automatically determine which one takes over the role of the access point~\cite{CampsMur2013,Strobel2024}.
However, Wi-Fi Direct has two main limitations: it does not support device discovery in the background if the utilizing application is stopped or not actively running in the foreground; another significant constraint of Wi-Fi Direct on modern mobile platforms is that it is not supported by iOS devices, as Apple developed a competing standard, AirDrop~\cite{airdrop}.
The Android equivalent of AirDrop, QuickShare, is developed cooperatively by Samsung and Google, and since November 2025 supports platform interoperability between Android and iOS, and has had support transfers between Androids, Windows computers, and Chromebooks before.
While similarly providing superior data rates, QuickShare, as it uses a combination of Wi-Fi Direct and Bluetooth, is still not usable for this work, as its core lies with Wi-Fi Direct, which requires more power and devices to have the using app constantly foregrounded.
Wi-Fi Aware (also known as Neighbor Awareness Networking / NAN) overcomes the background limitation of Wi-Fi Direct by utilizing a highly power-efficient \textit{discover-before-connect} paradigm.
Unlike Wi-Fi Direct, Wi-Fi Aware allows devices to automatically organize into synchronized local clusters and use a publish-and-subscribe mechanism to discover nearby services and exchange lightweight messages in the background, all without establishing a formal network connection.
Because it operates without maintaining an active connection state, it has a significantly lower power and memory footprint.
If devices need to share large amounts of data, they can seamlessly upgrade their discovery session into a fully bi-directional, high-throughput network connection.
Positioning itself as the most appropriate for our use-case, Wi-Fi Aware seems like a promising candidate; however, it was primarly intended for IoT devices~\cite{Saloni2016}, support for it on iOS was only recently introduced~\cite{WWDC25}, and thus it was deemed not widely supported enough.
Early research into ad hoc and delay-tolerant networks frequently relied on legacy 802.11b and 802.11g standards, which were highly capable for active data transfers but suffered from severe battery drain when devices were left idle~\cite{Ott2006}.
In later standard, 802.11n has introduced a significantly higher throughput and a 17\% improvement in energy efficiency over previous standards~\cite{Friedman2013}.

\subsection{Alternatives}
With Wi-Fi standards being exhaused for picking the base communication protocol, we evaluated other protocols such as Zigbee and LoRA, but very quickly forsaken due to lacking support for direct communication between mobile phones: Zigbee, LoRA, and 6LowPAN are all intended for small IoT devices to communicate among themselves or with centralized servers~\cite{Saloni2016}.

\subsection{Bluetooth}
A technology that is ubiquotous to all smartphones that have been released in the last 10 years is Bluetooth. Being standardized as Bluetooth 1.0 in 1999 by the Bluetooth Special Interest Group (SIG), Bluetooth was introduced to mobile phones by Ericsson in 2001 with the release of the T36, which allowed phones to replace bulky cables transfers with a wireless connection.
Since then, Bluetooth has seen four major version bumps:
\begin{itemize}
    \item Bluetooth 2.0 offered in comparison to Wi-Fi a much lower energy consumption for transferring large files, and higher resilience to signal interference from nearby transmissions.
    \item Bluetooth 3.0 was introduced as an upgraded specification designed to improve upon the throughput limitations of Bluetooth 2.x, doubling both send/receive throughput speeds to almost 200 KB/s, while the increasing the overall power consumption~\cite{Friedman2013}.
    \item Bluetooth 4.0 represented an elementary shift in Bluetooth technology, with the focus shifting to optimizing for ultra-low power consumption: Bluetooth Low Energy, or BLE, was introduced in this version. BLE presents until today the backbone for most modern smartphone p2p applications, due to major power-saving capabilities: because BLE supports background advertising and discovery without requiring apps to be actively foregrounded, BLE is the primary transport mechanism for decentralized networks and emergency systems like firechat, Briar, Corona-Warn-App, Twimight, BitChat, and more~\cite{firechat1,briar,CoronaWarnApp,bitchat}.
\end{itemize}

BLE is characterized by a dual-role split: advertising, which is carried out by the peripheral device, and scanning, which is done by the central device.
The central device controls the connection, while the peripheral provides the service.
Advertising happens on a set window, in which the peripheral sends out advertisements on an interval.
The advertisements contain a customizable device unique identifier, as well as custom manufacturer-specific data.
The BLE advertised services stem from the underlying General Attribute Protocol, or GATT.
The GATT protocol is a hierarchial protocol, where the highest level is GAP: General Attribute Profile. A GAP groups GATT services together, to offer a specific function, e.g. a heart rate monitor, or a grassroots networking profile.
Furthermore, GATT services group a set of characteristics, which are used for interfacing with data, and controlling settings pertaining to the service.
Each of the characteristic is identified by a UUID, a value, and has different read/write permissions; the data-flow between the peripheral and the central is bidirectional, with the peripheral being able to notify the central of any changes to a characteristic value as it changes, or central can request to read the current value of a characteristic, and more~\cite{BLE-guide}.
To see the peripheral's advertisements, the central device also has to scan during the same time: the scanning is also limited to a window, during which the central listens to all incoming advertisements from surrounding devices.
To discover GATT services, a device first must perform a connection to one such a device, which potentially includes pairing - depending on the permissions the device's GATT server defined.   
BLE presents advantages not only in its conservative use of battery power, but also in its security, and usage model. When an app registers a GATT serveer with the BLE adapter, the mobile phone controls it natively, and the app can remain backgrounded or on standby, while the OS continues to scan and advertise for services.
This allows for both a heightened security model due to the protocol abstraction by the operation system, as well as the grouping together of functionalities from different apps to a central native coordinator.
Providing a throughput of about 230 kbps at a low power usage~\cite{Jacopo2017}, BLE was able to answer all out criteria C1-C5.
Due to the aforementioned benefits of the protocol, we thus chose BLE as the application's main communication protocol.
\newcommand{\pmark}{$\sim$}    % partial / conditional

\section{Remote Peer-to-peer Communication Technologies}
\label{sec:remote-p2p-protocols}
After discussing the choice for a proximate communication protocol, this section continues with elaborating on how different protocols were evaluated to facilitate long-distance communication.
The protocols that were evaluated in this section include UDP, TCP, HTTPS, QUIC, and UDX.

\subsection{Application Protocols: HTTP, HTTPS, QUIC}
HTTP and HTTPS are the first protocols that we evaluated, as they lie at the base of every modern web applications, and have been tried-and-tested by many applications.
Most web servers started by using HTTP, and in a standard web application scenario with a typical client and server having continuous end-to-end connectivity it would be optimal; however to deploy it in a DTN environment where connectivity is intermittent at best, multiple auxillary systems would need to be used in order to provide hop-by-hop connectivity. As our goal is to provide connectivity with the least amount of external infrastructure, HTTP and its secure variant HTTP over TLS/TLS (HTTPS) require many subsystems that cannot be guaranteed to be freely available in our scenarios.
Additionally, HTTPS requires having one constantly listening TCP socket, plus one live TLS connection per peer, each with crypto state. Android and iOS tend to aggresively kill backgrounded applications and sockets. Recovering the socket connections would cost a whole TCP handshake, and a TLS handshake, which would cost multiple round trips per peer.
This shows that recovery for TCP-and-TLS-based protocols is a heavy operation that is very expernsive in terms of time-to-first-message after start.
QUIC, the transport standardized underneath HTTP/3, was designed to overcome two structural limitations of the TCP+TLS stack: its serialized handshakes and head-of-line blocking.
Built on top of UDP, QUIC folds the transport and cryptographic handshakes into a single exchange, establishing an encrypted connection in one round trip, and in zero round trips when resuming a previously cached session; it further multiplexes independent streams over one connection, so that a lost packet on one stream does not stall the others~\cite{Langley2017}.
However, QUIC remains connection-oriented: it maintains per-connection state, negotiates a full TLS~1.3 handshake on every cold start, and, like TCP, must re-establish that state whenever the underlying socket is torn down.
In a mobile setting where the operating system aggressively reclaims backgrounded sockets, this coupling of transport and cryptographic state reintroduces the very time-to-first-message penalty we sought to avoid.
The first protocol that we consideres for that, which already implements all of the above, is QUIC, or Quick UDP Network Connections~/cite{Langley2017}.
QUIC was designed with a mobile user in mind, supporting seamless network migration through the usage of Connection IDs to route and identify connections, rather than a source/destination IP:port + protocol 5-tuple combination; multiplexing multiple lightweight streams within a single connection allows to avoid the Head-of-Line Blocking issue, which causes queues of packets to buffer and be held up until the first packet is fully processed, severely limiting the performance, and is further exacerbated through out-of-order delivery; finally, in an effort to make common pairings more efficient, QUIC combines the transport connection establishment handshake with a cryptographic handshake to offer 0-RTT connectino setup, allowing data to be sent directly on the first message to known peers.
Since QUIC combines the cryptographic layer (normally TLS) with the network layer, it requires continuous cryptographic processing, which reportedly cost up to 3.5 times more CPU than normal TCP using TLS~\cite{Langley2017}. This computational exertion quickly drains mobile devices, since they are highly battery-constrained, and their CPU is much weaker than the desktop computers QUIC was designed for. Additionally, one of the main advantages of QUIC is the 0-RTT handshake data delivery, which depends on cached address tokens. These cached address tokens often become invalidated when mobile phones roam, dropping the success rate significantly~\cite{Langley2017}.
However, precisely because it combines so many things that UDP normally does not offer, QUIC became a monolith of a protocol that is mostly designed for stable end-to-end connections such as traditional HTTP (HTTP/3) client-server connections. 

\subsection{Transport-level Protocols: TCP, UDP}
The elementary transport protocols that enable QUIC and HTTPS are UDP and TCP respectively: to decouple cryptography from connection maintenance and application state, we decided to focus on the transport protocols that enable application communication.
Providing congestion-controlled, reliable, ordered, error-checked streams, TCP, the core transport protocol behind the IP protocol stack, was the first protocol that we investigated.
As a congestion-controlled transport protocol, TCP assumes that all transmission errors in its delivery path stem from link congestion in the network; that means, that for intermittent connectivity outages, TCP will engage in congestion-avoidance measures, which dramatically reduce the possible data rate on the link temporaily~\cite{RFC5681}.
Moreover, increasing the number of end-to-end hops in the network connection yields rapidly deterioratating performance with the increase of hops~\cite{Ott2010}.
Additionally, the usage of Network Address Translation (NAT) in order to assign multiple hosts to the same IPv4 address in a network makes it even more difficult to establish a p2p TCP connection without using any external proxies, as the NAT mapping is dependent on a 5-tuple: both peers IP addresses, ports, and protocol~\cite{Keidar2026}.
Because TCP maintains a strict connection state, the connection will break entirely any element in the 5-tuple formed by the underlying network changes.
Every time a mobile device roams, such as switching from a Wi-Fi network to cellular data, or changing its network attachment point, its IP address changes, forcing the TCP connection to be completely re-established from scratch.
To overcome the obstacles caused by node mobility and the ensuing change in Internet addresses, we have decided to focus on the UDP protocol as we expect frequent address changes due to our mobility model.

\subsection{UDX}
\label{sec:udx}
The transport-layer protocol that was chosen for the system, UDP, is a lightweight, data-oriented network protocol that allows the users to send datagrams without establishing a session or connection between each other. It prioritizes simplicity, incarnate in a lightweight header size, over error-correcting code, reliability, or packet-ordering which are all present in TCP, which makes it ideal for sending lightweight messages as it wastes no time establishing a reliable connection as its counterpart TCP.
UDP or UDP-based protocols are often used in time-critical applications like online gaming, streaming, and other p2p applications.
Due to the focus on simplicity, the resulting trade-off yields a transport-layer that, had it based itself directly on UDP, would have needed to implement error-correction codes, deal with network-level packet-loss, congestion control and flow control on active links. In addition to the transport-level control mechanisms, an encryption layer needs to be implemented on-top, in order to guarantee no adversary could read messages that were not intended for them.
The transport we ultimately selected for remote peer-to-peer links is UDX, the reliable stream transport that underpins the Holepunch/Pears networking stack~\cite{udx}.
UDX provides reliable, in-order, multiplexed, and congestion-controlled streams over plain UDP, much like QUIC, but with one deliberate design difference: it performs no handshake and no encryption of its own.
As the reference implementation phrases it, UDX has `no handshakes, no encryption, no features'~\cite{udx}, which is precisely what makes it inexpensive to (re-)establish.
Because UDX is connectionless, a stream is set up by each side independently creating a stream with a locally chosen identifier and connecting to the remote peer using the stream identifier and the remote socket address. Neither peer waits for an acknowledgement before it may send.
The very first packet can therefore already carry application payload, making the transport itself effectively \emph{zero round trip}.
Encryption and authentication are provided by separate layers: the Noise suite-based~\cite{noise} secure public key exchanged is performed after the application-level identity exchange happens, with session keys derived from the identity public keys through a birational map~\cite{Zafar2025}.
Crucially, this cryptographic handshake is not paid for directly in the transport layer, which makes it very simple to establish such a connection.

\section{Opportunistic Networking Systems}
\label{sec:oppnet}
\subsection{Background}
From the communication protocols we elaborate to research other existing opportunistic networking systems.
MANETs, who predate OppNets, seemed promising to enable ad-hoc communcication between mobile devices, requiring devices to forward messages for others in addition to allowing the devices themselves to communicate.
Talk about different routing strategies in MANETs?
One of the leading examples of MANETs that are still developed to this day, is the BATMAN-adv protocol (Better Approach To Mobile Ad-hoc Networking), by the Freifunk Gesellschaft.
In BATMAN, nodes constantly broadcast a hello message to their peers, and peers that receive those messages propagate those further. Peers measure the proportion of hello messages that they receive from a node to estimate their link quality, and broadcast that to their peers as well. Peers then re-broadcast the propagation from the peer with the highest ratio link quality, and thus a routing infrastructure is built iteratively.
Opportunistic Networks (OppNets) are a subsystem of DTNs, which are an evolution of legacy MANET systems.
In MANET systems there has to be a fully-connected end-to-end path before a message can be sent.
OppNets are an instance of the Store-Carry-Forward paradigm, where the obstacle that prevents MANETs from being used, is used as an advantage - node mobility. In Pocket-Switched Networks, an alias to OppNets, peers frequently change their network attachment point due to physical movements, as precisely due to the reason alluded to by their name -- the participating devices sit in the users' pockets: their smartphones. In those networks, peers carry messages with them as they move until the opportunity arises to forward the message to another peer~\cite{Boldrini2014}.
In OppNets, the emphasis lies on the opportunistic establishment of links between peers. Every possible participating user is henceforth referred to as a peer in the OppNet, and the links between them establish opportunistically: that means a peer randomally passing another peer will try to establish a (secure) connection to it.
Because OppNets rely on the physical mobility patterns of humans to establish bridges between disconnected clusters (either due to network partitions, physical distance, or lacking infrastructure), OppNets excel in decentralization of communication in areas with intermittent connectivity, such as disaster zones, remote areas, or regions that suffer from state control of the communication infrastructure~\cite{Strobel2024,Alsulami2022,Pozza2015}. The usage of mobility models to propagate information in a network is especially useful to establish information flows in regions that lack communication infrastructure, such as war zones, and disaster-struck areas.

\subsection{Routing in MANETs and OppNets}
\label{sec:opp-routing}
In MANETs, routing algorithms assume that the network topology formed is a single connected component, where all nodes can be reached, and strive to establish and maintain an efficient end-to-end path between source and destination by using control packets to evaluate the route~\cite{Ott2010,Almeida2012}.
MANET routing algorithms are usually categorized to either on-demand reactive algorithms, or proactive algorithms where nodes frequently send messages to their neighbors to evaluate their link quality.
The most common MANET routing algorithms are listed below:\begin{itemize}
\item \emph{AODV}: the Ad-hoc On-Demand Distance Vector algirithm is a \emph{reactive} protocol that discovers routes upon request of a node to send a message to a peer. It floods a Route Request message (RREQ) across the network: every node that forwards the RREQ add who they received it from, in order to record a reverse path. When the destination receives the first RREQ, it sends a Route Response (RREP) along the reverse path, and nodes that forward it back create an entry in their routing table, so that future requests could follow that path. To maintain the route, nodes that detect failures send back a Route Error (RERR) message, triggering another RREQ sequence~\cite{CMB-MANET}.
\item \emph{OLSR}: the Optimized Link State Routing algorithm is a \emph{proactive} protocol where all nodes create and maintain routing tables according to their own link quality, and focuses on evaluating the closest neighborhood for every node, making use of the notion that not every node needs to flood everything to everyone to build efficient routing tables. OLSR uses a Multipoint Relay (MPR) set of neighboring nodes that are well-connected to every two-hop neighbor, in order to form an efficient flooding backbone across the entire network.
\item \emph{Geographic}: use GPS location to restrict flooding done by routing protocols, for example by using building geolocation information and solving the routing problem for buildings~\cite{Liu2026}.
\end{itemize}
On the other hand, OppNet routing does not make the same assumptions: in OppNet routing, there is no end-to-end path that exists. Whether due to sparse connectivity graphs, intermittent network outages, or high node mobility, there is no instantenous end-to-path that connects all nodes, and the graph is comprised out of multiple components. Instead of dropping a packet when a path is not available to it, OppNet routing algorithms store those packets at a relaying node, the node carries it and forwards it opportunistically when encountering new peers.
There exist a plethora of OppNet routing algorithms, and we will cover the three most famous ones:\begin{itemize}
    \item \emph{Epidemic}: A flooding approach where a node floods copies of its messages to every other node it encounters. Achieving high delivery rates by spreading data like an epidemic disease, it is highly prone to congestion due to the massive amount of bandwidth and buffer space needed for an all-to-all broadcast~\cite{Ott2006,Almeida2012}.
    \item \emph{Spray-and-Wait}: This algorithm limits the number of copies of messages created, by creating a constant amount of copies for every message at an originator. In the first phase of the algorithm, the originator `sprays' copies upon encountering other peers: in every encounter, the originator hands over a half of all the copies it has to every peer, until exactly one copy remains, which shall be only forwarded directly to the destination. The relaying nodes continue with the same methodology, until the final node is reached.
    \item \emph{PRoPHET} (Probabilistic Routing Protocol using History of Encounters and Transitivity): The PRoPHET protocol is a model-based routing algorithm, which uses stochastics to improve the probability of successful delivery, by recording past contact history between nodes and using nodes with the highest probability of eventually reaching a destination efficiently through their mobility pattern.
\end{itemize}
In this thesis, we use epidemic routing because it has the highest likelihood to deliver a message, at the cost of network congestion, bandwidth, and buffer space due to the DTN stores at every node buffering messages. That is due to focusing on broadcast dissemination, as our goal is to ensure every message is delivered to every unique node in a network at least once~\cite{Schmidheiser2025}. Cases where such a dissemination model is useful is in emergency situations, where disaster alerts have to be spread to as many people as possible as fast as possible, and protest scenarios where it must be ensured that all protestors receive vital information, such as evidence of police brutality, and coordination efforts.
Additionally, to create a research baseline, using an epidemic flooding model is beneficial because it lets researchers compare more advanced routing algorithms to a baseline naïve model, which acts as an upper bound for delivery rates and network congestion.
Early examples of OppNets have leveraged the notion of Store-Carry-Forward to offer intermittent connectivity without central communication infrastructures~\cite{Juang2002}, providing wide-scale wildlife tracking during mass migrations~\cite{Pozza2015}, as well as voicemail and email services to rural areas in developing countries~\cite{Lindgren2007,Pentland2004}.
On smartphones, one of the most famous and large-scale deployments was Firechat~\cite{firechat1}, which was broadly used in the 2014 Hong Kong protests, and later in Iran and Iraq.
Firechat allowed users to use either direct p2p communication, broadcast a message to all mesh participants, or use IRC-style chatrooms based on proximity. Using Firechat, tens-of-thousands of protestors were able to organize and coordinate supplies for the dissidents during the protests, during a time of Internet blackout scare by the government, as protestors chanted for democracy in the Occupy Central protests.
However, FireChat was not meant for this use case, as it was lacking basic encryption, privacy, and resilience~\cite{firechat2}.
From the technical perspective, FireChat utilized similar WiFi and Bluetooth to forward messages over multiple hops without cellular service using a store-and-forward mechanism, using Apple's Multipeer Connectivity framework.
Focusing on the ensuring free information flow in protest use-case, Bridgefy used a similar tech stack but did support encryption: it was used extensively again in the 2019-2020 Hong Kong protests and during Citizenship Amendment Act protests in India 2019.
As research shows~\cite{Albrecht2021}, adversaries were able to completely shut Bridgefy down using a single well-crafted message delivered across the mesh, deanonymize users, reconstruct users' social graphs, decrypt and read messages.
Other notable academic and practical systems include the Serval Project's BatPhone~\cite{batphone}, which enables infrastructure-less VoIP mobile telephony for disaster relief using the BATMAN protocol, and Twimight, an open-source Twitter client featuring a `disaster mode' that allows tweets to spread epidemically from phone to phone via Bluetooth when traditional networks fail~\cite{Legendre2011}.
Building upon years of experience as the CEO of Twitter, Jack Dorsey, together with permissionlesstech, released BitChat in July 2015, a highly secure evolution of the aforementioned platforms. BitChat features a dual-stack transport over both a BLE mesh and the WiFi medium, using decentralizede Nostr relays to encrypt and anonymize messages when the user is Internet access is available.
To efficiently flood messages across the local mesh without overwhelming the network, BitChat implements an epidemic-like routing protocol, which widely distributes data while using node-specific Bloom filters to quickly identify and reject duplicate packets.
Furthermore, by natively integrating robust end-to-end encryption and mutual authentication via the Noise Protocol Framework and Noise key exchange, BitChat solves the glaring privacy vulnerabilities that plagued earlier mesh apps.
This advanced, censorship-resistant architecture has already seen significant real-world adoption, proving vital for coordination during large-scale protests in Nepal and state-ordered internet blackouts in Iran and Uganda.

\section{Decentralized Social Networks}
\label{sec:dsn}
Traditional Centralized Social Networks (CSNs) are highly subsceptible to Internet outages, control of the underlying infrastructure to monitoring, surveillance, and censorships due to being regulated or controlled by government organizations. Because all data flows through centralized servers in CSNs, authoritan regimes can easily force companies by restricting financial access, or order executive shutdowns at the network level~\cite{Lopez2026}. WHen communication is halted in this manner, human social movements are distrupted, and the fundamental right to freedom of (digital) assembly is eliminated without a visible act of physical repression. Governmental control oftentimes leads to repression of human rights during times of political unrest, as access to such information gives governments more power to crack down on political activists and the general population~\cite{digital-repression-palestine,internet-shutdown-human-rights}. CSNs require a large set of infrastructure to support the communication backbone that enables the population access to them: more than 3 billion monthly active users use Facebook~\cite{FBMAU}, and access to the aforementioned infrastructure oftentimes requires governmental cooperation through facilitation with network authorities or ISPs.
Twitter (now X), after being acquired by Elon Musk, has underwent a massive change in its content policy, moderation, and workforce --- which lead to millions of users who left the platform~\cite{Balduf2024}. As a centralized private online platform, Twitter (and others) are sensitive to (hostile) takeovers that can leave users feeling that the platform no longer represents their values.
Despite reaching unprecedented global reach, CSNs are constrained by their fundamental dependency on a single controling authority, which gives rise to severe societal and privacy challenges. Most notably, the free-to-use, user-as-a-product business model that traditional CSNs use, make use of user data in order to market personalized ads to unsuspecting users, by using embedded tracking and analytic software.
This business model, coined Surveillance Capitalism, exploits user data, online behavior, and original content to provide financial value for the platform, with little to no renumeration for the creators, as users are forced to surrender ownership of their data to third-party operators~\cite{Zuboff2023}.
Additionally, the platform operators wield the controlling power of information, by providing algorithmic curation and content moderation, thereby creating opaque governance that gives rise to environments which are vulerable to ideological polarization, in which users have limited agency over their digital lives~\cite{Balduf2024,Jeong2025}.
To counteract the aforementioned vulnerabilities and business practices, Decentralized Social Networks (DSNs) have evolved to distribute data and control across a variety of different entities, to restore user autonomy, data sovereignty, and improve censorship resistance. Based on recent literature~\cite{Jeong2025}, DSNs can be categorized into four primary architectural paradigms, each with their own technical and societal trade-offs, and a minimum set of essential agents (the minimal number of nodes whose removal paralyzes the network)~\cite{Shapiro2026}. These platforms diverge in their used data structures, content moderation policies, centralization features, and reliance on infrastructure.
The first category is the federated architecture, in which networks distribute control over multiple independently operated servers that communicate over standardized protocols.
In the federated architecture, control is held by local operators of the platform, which corresponds to an unlimited, but not universal, set of operating agents~\cite{Shapiro2026}.
A popular example of this architecture is Mastodon, which is comprised out of distributed home servers which all run the same Mastodon code, that interoperate using a common protocol. % , allowing users that register with a home server to follow users both from their home server, as well as from a different server.
Subverting the centralization in common platforms like Facebook, X, and Instagram, Mastodon distributes control across multiple independently operated servers.
Users on Mastodon register with a home server to store their profile, settings, and post data, and can follow users from all servers. Following a user on another instance relies on the home server performing the subscription on their behalf, which requires the server admins to enable federation, and the remote instance to not be blocked by the home instance's instance blocking feature~\cite{Aravindh2019}.
Thus, users on a Mastodon server are governed by an administrator (or a set thereof) that set basic terms of usage, user settings, and form a sort of constitution for using them (i.e. terms of agreement).
While preventing a single server from owning the entire network, empirical measurements of Mastodon reveal that the user population is largely skewed: just 10\% of instances host almost half of all users, and the top 5\% of instances generate nearly 95\% of all posts~\cite{Aravindh2019}. Additionally, the uncoordinated, voluntary nature of the instance administrators lead to instances being heavily co-located within a small set of Autonomous Systems (ASes), with AWS alone accounting for 30\% of all global users: this created massive single points of failure, with an outage affecting just the top 10 of global instances wiping out over 62\% of all global posts~\cite{Aravindh2019}.
The empirical data highlights a gap in federated architecture models: without global decentralized mechanisms that synchronize state (i.e. content replication) and perform load balancing across instances, federated isolated networks remain highly sensitive to data loss caused by infrastructure outages.
The second category of DSNs is the peer-to-peer (p2p) architecture, which operates without servers, relaying on its users' end devices to form a distributed overlay. This category of DSNs has an infinite, universal amount of essential agents, as the network is completely decentralized~\cite{Shapiro2026}. In the p2p architecture, data is stored locally on user devices, and is relayed to peers directly upon rendezvous. Examples of the p2p architecture include the Secure Scuttlebutt (SSB) protocol, Nostr, and Briar: the foundation of SSB is secure append-only logs that are synchronized between peers (oftentimes over TCP, sometimes over other protocols i.e. Bluetooth Classic in Manyverse) in a manner similar to `personal blockchains'.
SSB is connection based, and requires users to actively block users they do not wish to communicate with, since it assumes all-to-all communication: unless limited, the user's log updates are conferred to all other users upon rendezvous. In comparison, Nostr uses a set of global relays that forward cryptographically signed lightweight messages without maintaining global state, while Briar uses WiFi and Bluetooth Classic for close range interactions, and the Tor decentralized network of onion routers to secure relay information over the Internet~\cite{briar}.
Using online relays or lightweight servers to forward messages, platforms like Nostr and Briar suffer from a critical dependence on third party providers, which renders them mostly useless in Internet outages.
Without a centralized `permissioned' administration, optimizing for maximum user autonomy, privacy, and censorship resistance, the p2p network makes content availability depend strictly on peers being online; creating a global aggregated view of all posts to discover community trends or enforce network-wide moderation is impossible~\cite{Jeong2025}.
The third architecture of DSNs is the blockchain architecture, in which user interactions like posts and likes are recorded as immutable transactions on a public distributed ledger, verified by miners or validators, and ensuring global consistency and transparency~\cite{Jeong2025}.
In the blockchain architecture, the number of essential agents is constrained and is finite, and largely corresponds to the bootstrap nodes~\cite{Shapiro2026}.
Such blockchain systems inherently solve the problem of global trend discovery and data availability, ensuring data stays available forever, but face a clash with rights such as the right to be forgotten, as due to the chain's immutability harmful content cannot be deleted once created. Additionally, platforms that use plutocratic, token-based incentives (e.g. Steemit) are often controlled or manipulated by `whale' users: users with a large holding of tokens, that allows them to maximalize their posts visibility and reachability~\cite{Jeong2025}. That creates a situation in which rich users (possessing a high amount of `coins') have a higher leverage in the platform, trumping poor users. The ability to purchase the platform coins with traditional currency sparks the notion of a plutocracy, in which great wealth dictates control of the government and public discourse~\cite{Shapiro2026}.
Finally, the fourth architecture of DSNs is the hybrid architecture: platforms belonging to this category attempt to balance the scalability of centralized systems with the user autonomy of decentralized protocols. Bluesky is a major proponent of this architecture: in Bluesky, the different network components are decomposed into independent sub-components that interact. In Bluesky, user identity is secured via Decentralized Identifiers (DIDs), and can be securely transferred between different Personal Data Servers (PDses); using the AT protocol, user data is heaved above residence at traditional data servers, and is located at the user's PDS provider, which Bluesky encourages users to run themselves. Using the AT protocol, the user can decide to allow Bluesky access to their data identified by their DID record, which allows a centralized component, the Firehose, to efficiently aggregate and cache their data. Additionally, content curation and moderation are decentralized and decoupled: users can subscribe to `Labeler' and `Feed Generators' accounts, which respectively customize their safety filters and feed generating algorithms, by pulling post labels from third party accounts (and manually deciding what to do with each label generated), and curating content from accounts the user follows~\cite{Balduf2024}.
An additional, a rather new proponent of this architecture is BitChat, which is optimized for censorship resistance using a dual-transport protocol. When the Internet is available, BitChat routes messages globally using the aforementioned Nostr relay network, and falls back to offline BLE mesh~\cite{Biagioni2026}. In the latter mode, BitChat uses an Epidemic-like protocol (`Efficient Gossip with Bloom Filters') in order to propagates messages to remote peers, ensuring vital communication persists despite unavailability of local Internet infrastructure.

\section{Grassroots Social Networks}
\label{sec:gsn}
To overcome CSNs, the congestion encumbering mesh flooding-based DSNs, the plutocracy of blockchain-based social networks, and the relay dependency of p2p DSNs, Grassroots Social Networks (GSN) proposes a lean p2p architecture, without relaying on any global infrasatructure, third-party relays, or shared global ledgers. GSNs operating strictly on the smartphones of its users, granting total user agency over personal information and social graphs, and thus has a universal, infinite amount of essential agents~\cite{Shapiro2026}.
Grassroots Social Networks aim to provide an egalitarian, cooperative, and democratic alternative to the aforementioned DSNs and CSNs. Unlike in other platforms, by using the GSN any user agent initiates their own instance, choosing their own cryptographic key pair, letting agents operate independently and in isolation of each other by will, sharing data by choice. GSNs are informally defined by having an agreed constitution (the set of basic governing rules that dictate how agents operate): due to the constitution, each user can participate in an endless amount of independent instances, always producing a valid set of operations and state mutations, so that a coalescing of users into a group always produces a seamless, valid `global' state.
This work positions itself as networking layer for a larger Grassroots Social Network, which is currently being researched with Ehud Shapiro (Weizmann Institute, Israel and London School of Echonomics, UK), Idit Keidar (Technion - Israel Institute of Technology), and Ohad Eitan (Technion - Israel Institute of Technology).
One of the bigger challenges in both p2p and grassroots social networks, is devising a key management protocol which is secure against major faults~\cite{Inyangson2025}: the loss of a smartphone (state loss), or a private cryptographic key (identity loss) - both of which would typically result in the permanent loss of the user's digital identity and complete social graph. 
In~\cite{Eitan2026}, Eitan et al. devise a solution to the problem of recovering an account with no available central server to reset a password or restore data. By building on the user's social graph, identity can be recovered by reaching a supermajority of `Identity Custodians': a designated set of trusted friends specified by the user, that authorize a newly-created public key as the reincarnation of an old, discarded key, preserving all of the user's friendships. If a user still has their private key, state can be recovered using the user's friends as `State Custodians', which uses them to recreate the user's social graph and data using the user's private key.
Concretely, this work provides an overview of the Grassroots Networking Layer (GSL) built to support GSN: a hybrid, dual-stack system supporting two communication protocols. In close proximity, a BLE mesh is used to bootstrap friendships and exchange messages containing text and media. When Internet is available, direct IP communication over UDP is supported, after hole-punching is facilitated using connected friends or BLE as a signaling channel.